% Source PDF: source/普通高考/2013/2013天津理.pdf, page 1 (question 3).
% The source PDF contains vector line art (pdfimages -list returned no image objects).
% Node -> directed-edge list: 开始 -> 输入 x -> S=0 -> S=S+x^3 -> S>=50?;
% 是 -> 输出 S -> 结束; 否 -> x=2x -> loop-back to the stem immediately above
% S=S+x^3.  The initial S=0 edge and the loop-back share this stem but carry
% distinct directed segments, matching the source topology.
% solid segments: all node outlines and connectors; dashed segments: none.
% labels: 是 beside the downward decision exit; 否 beside the rightward exit;
% all node text is locally zero-indented and centered.

\begin{tikzpicture}[
  x=1cm,
  y=1cm,
  >=Stealth,
  line width=0.45pt,
  line cap=round,
  line join=round,
  font=\normalsize,
  flowtext/.style={anchor=center, align=center,
    execute at begin node={\setlength{\parindent}{0pt}}},
  flowbox/.style={draw, minimum height=0.72cm, inner xsep=4pt, inner ysep=2pt,
    align=center, execute at begin node={\setlength{\parindent}{0pt}}},
  terminal/.style={flowbox, rounded corners=0.18cm, minimum width=1.30cm},
  io/.style={flowbox, trapezium, trapezium left angle=72, trapezium right angle=108,
    minimum height=0.78cm},
  decision/.style={flowbox, diamond, aspect=2.75, minimum width=3.65cm,
    minimum height=1.05cm, inner xsep=1pt, inner ysep=1pt},
  branchlabel/.style={flowtext, inner sep=0pt},
]
  % Main sequence: the coordinates preserve the source's centered vertical stem.
  \node[terminal] (start) at (0,15.0) {开始};
  \node[io, minimum width=2.00cm] (input) at (0,13.45) {输入 $x$};
  \node[flowbox, minimum width=2.45cm] (init) at (0,11.85) {$S=0$};
  \node[flowbox, minimum width=3.70cm] (update) at (0,9.45) {$S=S+x^3$};
  \node[decision] (test) at (0,7.35) {$S\geq 50?$};
  \node[io, minimum width=2.00cm] (output) at (0,4.05) {输出 $S$};
  \node[terminal] (finish) at (0,2.15) {结束};

  \node[flowbox, minimum width=2.60cm] (double) at (4.35,9.45) {$x=2x$};

  \draw[->] (start.south) -- (input.north);
  \draw[->] (input.south) -- (init.north);

  % S=0 enters the shared stem; loop-back returns to the same merge point.
  \coordinate (merge) at (0,10.55);
  \draw (init.south) -- (merge);
  \draw[->] (merge) -- (update.north);

  \draw[->] (update.south) -- (test.north);

  % Yes branch: decision -> output -> finish.
  \draw[->] (test.south) -- (output.north);
  \draw[->] (output.south) -- (finish.north);

  % No branch: decision -> x=2x -> leftward loop-back into the shared stem.
  \draw[->] (test.east) -- (4.35,7.35) -- (4.35,8.95) -- (double.south);
  % The loopback arrow terminates at the shared stem.
  \draw[->] (double.north) -- (4.35,10.55) -- (-0.35,10.55) -- (merge);

  \node[branchlabel, anchor=west] at (0.42,5.85) {是};
  \node[branchlabel, anchor=west] at (2.70,7.78) {否};
\end{tikzpicture}
